\documentclass{article}
\usepackage{graphicx} % Required for inserting images
\usepackage{amsmath}
\usepackage{booktabs}

\title{Process intensification of catalytic-wall Taylor-Couette reactors by reinforcement learning}
\author{Isaiah Bascom}
\date{July 2026}

\begin{document}

\maketitle

\section{Abstract}
Previous work intensified reactions in catalytic-wall Taylor-Couette reactors by modulating the angular velocity of the reactor's inner cylinder, saving a considerable amount of energy~\cite{lopez2024process}.
Here we extend their methods by creating a RL system that chooses unconventional modulation strategies that requires less energy than previous benchmark policies.

\section{Introduction}
Process intensification tends to be an expensive endeavor implying replacing old equipment with new and improved ones. 
Here we offer a treat giving you similar scale of improvement without asking for your company credit score.
Taylor-Couette reactors are super predictable and consistent so they have become very popular in industrial processes such as crystallization, electrolysis, and catalysis. 

\newpage
\section{Methodology}
\subsection{Problem Formulation}
    For the sake of knowing whether our RL approach outperforms Lopez-Guajardo et al.~\cite{lopez2024process}, 
    our formulation closely follows theirs. The Taylor-Couette reactor
    geometry is summarized in Table~\ref{tab:geometry}.

\begin{table}[ht]
    \centering
    \caption{Geometry of the catalytic-wall Taylor-Couette reactor. The
    computational domain is a $5^\circ$ azimuthal wedge of the full annulus,
    with a one-cell side-outlet band at the bottom of the outer wall.\newline}\label{tab:geometry}
    \begin{tabular}{lcc}
        \toprule
        Parameter & Symbol & Value \\
        \midrule
        Inner cylinder radius   & $R_i$          & $25.40$ mm \\
        Outer cylinder radius   & $R_o$          & $31.75$ mm \\
        Gap width               & $d = R_o - R_i$ & $6.35$ mm \\
        Reactor height          & $H$            & $38.10$ mm \\
        Radius ratio            & $\eta = R_i/R_o$ & $0.80$ \\
        Aspect ratio            & $\Gamma = H/d$ & $6.0$ \\
        Side-outlet band height & $h_\mathrm{out}$ & $0.42$ mm \\
        Azimuthal wedge angle   & $\theta$       & $5^\circ$ \\
        \bottomrule
    \end{tabular}
\end{table}

The incompressible Navier-Stokes equations
\begin{align}
    \rho \nabla \cdot \mathbf{u} &= 0, \\
    \rho \left( \frac{\partial \mathbf{u}}{\partial t} + \mathbf{u} \cdot \nabla \mathbf{u} \right) &= -\nabla p + \mu \nabla^2 \mathbf{u},
\end{align}

and the transport of dilute species equation
\begin{align}
    \frac{\partial c}{\partial t} + \mathbf{u} \cdot \nabla c &= D \nabla^2 c,
\end{align}

were solved by FVM with the OpenFOAM v2506~\cite{weller1998openfoam} framework.


The Reynolds number was defined as $Re = \frac{\omega R_i d}{\nu}$, 
where $\omega$ is the angular velocity of the inner cylinder and $\nu$ is the kinematic viscosity of the fluid. 
The flow was assumed to be laminar and incompressible, 
and the no-slip boundary condition was applied at both cylinder walls.

The Farama Foundation's Gymnasium~\cite{towers2024gymnasium}
was used as the framework for our reinforcement learning (RL) environment with 


\section{Results}

\section{Discussion}

\bibliographystyle{plain}
\bibliography{references}

\nocite{lopez2024process,fujimoto2018addressing,lee2023turbulence,zhou2025deep,bergman2011fundamentals,towers2024gymnasium,weller1998openfoam}


\end{document}